% Figura: fw-modulos -- requiere % =====================================================================
%  Estilos compartidos por los diagramas TikZ de los capitulos.
%  Va en el PREAMBULO del documento principal:
%      % =====================================================================
%  Estilos compartidos por los diagramas TikZ de los capitulos.
%  Va en el PREAMBULO del documento principal:
%      % =====================================================================
%  Estilos compartidos por los diagramas TikZ de los capitulos.
%  Va en el PREAMBULO del documento principal:
%      \input{diagramas/estilos-diagramas}
% =====================================================================

\usepackage{tikz}
\usepackage{amsmath}
\usepackage{amssymb}
\usetikzlibrary{arrows.meta,positioning,shapes.geometric,shapes.misc,calc,fit,backgrounds,decorations.pathreplacing}

\definecolor{cbloque}{HTML}{E8EEF4}
\definecolor{cborde}{HTML}{4A6D8C}
\definecolor{cacento}{HTML}{D9E6D5}
\definecolor{cacentob}{HTML}{5A7A4A}
\definecolor{cgris}{HTML}{EDEDED}
\definecolor{cgrisb}{HTML}{888888}
\definecolor{calarma}{HTML}{F6DEDE}
\definecolor{calarmab}{HTML}{A03030}

\tikzset{
  bloque/.style={draw=cborde,fill=cbloque,rounded corners=2pt,align=center,
                 inner sep=5pt,font=\footnotesize,minimum height=9mm},
  bloqueg/.style={bloque,draw=cgrisb,fill=cgris},
  bloquea/.style={bloque,draw=cacentob,fill=cacento},
  bloquer/.style={bloque,draw=calarmab,fill=calarma},
  proceso/.style={bloque,minimum width=52mm},
  decision/.style={draw=cborde,fill=cbloque,align=center,font=\scriptsize,
                   diamond,aspect=2.2,inner sep=1pt},
  flecha/.style={-{Stealth[length=2mm]},draw=cborde!80!black,thick},
  flechag/.style={flecha,draw=cgrisb},
  et/.style={font=\scriptsize,inner sep=1.5pt,fill=white,align=center},
  etl/.style={et,midway},
  grupo/.style={draw=cgrisb,dashed,rounded corners=3pt,inner sep=4mm},
  tit/.style={font=\scriptsize\bfseries,text=cgrisb!60!black},
}

% --- utilidades para diagramas de secuencia ---
\tikzset{
  vida/.style={draw=cgrisb,dash pattern=on 2pt off 2pt},
  cab/.style={bloque,minimum width=26mm,font=\scriptsize\bfseries},
  nota/.style={draw=cacentob,fill=cacento,font=\scriptsize,align=center,
               rounded corners=1pt,inner sep=3pt},
}
\newcommand{\msg}[4]{\draw[flecha] (#1,#3) -- (#2,#3) node[et,midway,above=0pt]{#4};}
\newcommand{\msgr}[4]{\draw[flecha,dashed] (#1,#3) -- (#2,#3) node[et,midway,above=0pt]{#4};}
\newcommand{\auto}[3]{%
  \draw[flecha] (#1,#2) -- ++(0.55,0) -- ++(0,-0.42) -- (#1,{#2-0.42});
  \node[et,anchor=west] at ({#1+0.7},{#2-0.21}) {#3};}


% =====================================================================

\usepackage{tikz}
\usepackage{amsmath}
\usepackage{amssymb}
\usetikzlibrary{arrows.meta,positioning,shapes.geometric,shapes.misc,calc,fit,backgrounds,decorations.pathreplacing}

\definecolor{cbloque}{HTML}{E8EEF4}
\definecolor{cborde}{HTML}{4A6D8C}
\definecolor{cacento}{HTML}{D9E6D5}
\definecolor{cacentob}{HTML}{5A7A4A}
\definecolor{cgris}{HTML}{EDEDED}
\definecolor{cgrisb}{HTML}{888888}
\definecolor{calarma}{HTML}{F6DEDE}
\definecolor{calarmab}{HTML}{A03030}

\tikzset{
  bloque/.style={draw=cborde,fill=cbloque,rounded corners=2pt,align=center,
                 inner sep=5pt,font=\footnotesize,minimum height=9mm},
  bloqueg/.style={bloque,draw=cgrisb,fill=cgris},
  bloquea/.style={bloque,draw=cacentob,fill=cacento},
  bloquer/.style={bloque,draw=calarmab,fill=calarma},
  proceso/.style={bloque,minimum width=52mm},
  decision/.style={draw=cborde,fill=cbloque,align=center,font=\scriptsize,
                   diamond,aspect=2.2,inner sep=1pt},
  flecha/.style={-{Stealth[length=2mm]},draw=cborde!80!black,thick},
  flechag/.style={flecha,draw=cgrisb},
  et/.style={font=\scriptsize,inner sep=1.5pt,fill=white,align=center},
  etl/.style={et,midway},
  grupo/.style={draw=cgrisb,dashed,rounded corners=3pt,inner sep=4mm},
  tit/.style={font=\scriptsize\bfseries,text=cgrisb!60!black},
}

% --- utilidades para diagramas de secuencia ---
\tikzset{
  vida/.style={draw=cgrisb,dash pattern=on 2pt off 2pt},
  cab/.style={bloque,minimum width=26mm,font=\scriptsize\bfseries},
  nota/.style={draw=cacentob,fill=cacento,font=\scriptsize,align=center,
               rounded corners=1pt,inner sep=3pt},
}
\newcommand{\msg}[4]{\draw[flecha] (#1,#3) -- (#2,#3) node[et,midway,above=0pt]{#4};}
\newcommand{\msgr}[4]{\draw[flecha,dashed] (#1,#3) -- (#2,#3) node[et,midway,above=0pt]{#4};}
\newcommand{\auto}[3]{%
  \draw[flecha] (#1,#2) -- ++(0.55,0) -- ++(0,-0.42) -- (#1,{#2-0.42});
  \node[et,anchor=west] at ({#1+0.7},{#2-0.21}) {#3};}


% =====================================================================

\usepackage{tikz}
\usepackage{amsmath}
\usepackage{amssymb}
\usetikzlibrary{arrows.meta,positioning,shapes.geometric,shapes.misc,calc,fit,backgrounds,decorations.pathreplacing}

\definecolor{cbloque}{HTML}{E8EEF4}
\definecolor{cborde}{HTML}{4A6D8C}
\definecolor{cacento}{HTML}{D9E6D5}
\definecolor{cacentob}{HTML}{5A7A4A}
\definecolor{cgris}{HTML}{EDEDED}
\definecolor{cgrisb}{HTML}{888888}
\definecolor{calarma}{HTML}{F6DEDE}
\definecolor{calarmab}{HTML}{A03030}

\tikzset{
  bloque/.style={draw=cborde,fill=cbloque,rounded corners=2pt,align=center,
                 inner sep=5pt,font=\footnotesize,minimum height=9mm},
  bloqueg/.style={bloque,draw=cgrisb,fill=cgris},
  bloquea/.style={bloque,draw=cacentob,fill=cacento},
  bloquer/.style={bloque,draw=calarmab,fill=calarma},
  proceso/.style={bloque,minimum width=52mm},
  decision/.style={draw=cborde,fill=cbloque,align=center,font=\scriptsize,
                   diamond,aspect=2.2,inner sep=1pt},
  flecha/.style={-{Stealth[length=2mm]},draw=cborde!80!black,thick},
  flechag/.style={flecha,draw=cgrisb},
  et/.style={font=\scriptsize,inner sep=1.5pt,fill=white,align=center},
  etl/.style={et,midway},
  grupo/.style={draw=cgrisb,dashed,rounded corners=3pt,inner sep=4mm},
  tit/.style={font=\scriptsize\bfseries,text=cgrisb!60!black},
}

% --- utilidades para diagramas de secuencia ---
\tikzset{
  vida/.style={draw=cgrisb,dash pattern=on 2pt off 2pt},
  cab/.style={bloque,minimum width=26mm,font=\scriptsize\bfseries},
  nota/.style={draw=cacentob,fill=cacento,font=\scriptsize,align=center,
               rounded corners=1pt,inner sep=3pt},
}
\newcommand{\msg}[4]{\draw[flecha] (#1,#3) -- (#2,#3) node[et,midway,above=0pt]{#4};}
\newcommand{\msgr}[4]{\draw[flecha,dashed] (#1,#3) -- (#2,#3) node[et,midway,above=0pt]{#4};}
\newcommand{\auto}[3]{%
  \draw[flecha] (#1,#2) -- ++(0.55,0) -- ++(0,-0.42) -- (#1,{#2-0.42});
  \node[et,anchor=west] at ({#1+0.7},{#2-0.21}) {#3};}

 en el preambulo
\begin{tikzpicture}[x=1cm,y=1cm]

\node[bloque,text width=52mm] (main) at (0,3.2)
  {\texttt{hsrl-control.c}\\[1pt]\scriptsize \texttt{main()}, lazo de ciclo, parser de
   comandos en \texttt{core1}, línea \texttt{[log]}};

\node[bloque,text width=32mm] (adc)  at (-4.2,1.0) {\texttt{adc\_capture.c/.h}\\[1pt]\scriptsize captura de picos\\sincronizada al trigger};
\node[bloque,text width=32mm] (ctrl) at (0,1.0)    {\texttt{control.c/.h}\\[1pt]\scriptsize máquina de estados\\y cálculo de deltas};
\node[bloque,text width=32mm] (seed) at (4.2,1.0)  {\texttt{seeder\_comm.c/.h}\\[1pt]\scriptsize protocolo SCPI\\sobre UART};

\node[bloquea,text width=44mm] (cfg) at (-2.4,-1.2)
  {\texttt{config.h}\\[1pt]\scriptsize pines, parámetros por defecto y límites};
\node[bloqueg,text width=44mm] (sdk) at (3.2,-1.2)
  {\textbf{Pico SDK 2.3.0}\\[1pt]\scriptsize \texttt{hardware\_adc}, \texttt{hardware\_uart},\\
   \texttt{hardware\_sync}, \texttt{pico\_multicore}};

\draw[flecha] (main) -- (adc);
\draw[flecha] (main) -- (ctrl);
\draw[flecha] (main) -- (seed);
\draw[flechag] (adc)  -- (cfg);
\draw[flechag] (ctrl) -- (cfg);
\draw[flechag] (seed.south) -- (sdk.north);
\draw[flechag] (adc.south) -- ++(0,-0.55) -| (sdk.150);
% rodea el diagrama por la izquierda: pasaba por dentro de adc y de config.h
\draw[flechag] (main.west) -- (-6.9,3.2) -- (-6.9,-1.2) -- (cfg.west);


\end{tikzpicture}